\scanpage{295}
\mspage{19}
% archive sheet number: 294

On trouve, quand on dispose de $\mathrm{Loc}(6,a)$\inserted{(\unsure{sa forme|ses formes} \unsure{duale|duales}\inserted{i.e. pour les \uline{fermés}}, c'est kif-kif)}, que les \uline{immersions ouvertes de Dwyer} sont dans $\mathrm{Cof}_{\unsure{W|w}}$, \deleted{\unsure{d'où}} ce qui suffit pour avoir \uline{une \inserted{bonne} théorie} \uline{des cofibres} $W$-\uline{homotopiques}, \uline{dans \unsure{laquelle|lesquelles} les carrés} $\mathrm{MV}$ \uline{sont} $W$-\uline{cocartésiens}. (chose à laquelle on tient absolument).

Quant à $\mathrm{Loc}(6b)$, \deleted{\unsure{c'est}} si je dispose à la place de $\mathrm{Loc}(6'b)$, \inserted{\uline{et $\mathrm{Loc}(6a)$,}} je dois m'en tirer pour prouver $\mathrm{Loc}(6)$
\deleted{et \unsure{même}, pour \unsure{avoir} $\mathrm{Loc}(6)$ sous la forme en apparence plus forte de l'ancien axiome $W(7)$ (``axiome du cube'' ou ``lemme des cinq'', cf \uline{XII} p.~76, \unsure{i.e.|ici} aussi la forme}
\[
\deleted{\text{c)}\quad f_0,f_3\in W\Longrightarrow f_1,f_2\in W.}
\]

Dans XII Appendice (pages \unsure{265 ff|265 A}) je prouve, paraît-il, que (avec les notations actuelles) $\mathrm{Loc}(6')$ implique $\mathrm{Loc}(6)$, et $\unsure{=|\Longleftrightarrow}\mathrm{Loc}(6c)$ si $W$ stable par facteurs directs (condition $\mathrm{Loc}(1\,\mathrm{bis})$). Mais la démonstration est incomplète, puisque j'admets tacitement que les immersions de Dwyer sont des $\mathrm{Cof}_W$, chose que je n'avais établie que moyennant l'ax. $W(7)$, en fait moyennant $\mathrm{Loc}(6'a)$, ce qui est une partie de $W(7)$).
